% Where the theory stands (written 2026-10-08, after the head/tail correction)
\section{Where the theory stands (2026-10-08)}\label{sec:status}
\paragraph{The sign issue in one paragraph.} The defect angle in \texttt{defects.dat} is the direction of $\nabla\!\cdot Q$, which for a $+\half$ points
to its \emph{tail} (the splay side). Verified three ways: the formula in the code, the director field around bulk $+\half$ (director along the axis on the
recorded side, $|\mathbf n\cdot\hat e|=0.92$--0.95, across it on the other, 0.20), and the motion of bulk $+\half$ (against the recorded angle, 66--72\% of steps,
$\approx60\sigma$). Head, self-propulsion and active force point along $-\hat e$. Everything that only uses positions, charges or relative angles is
unaffected (distributions, charge, survival, depth law, shielding, fluctuations, the tilt--position lock); everything that turned a $+\half$ into a force
had the wrong sign (the point-force Stokes models, the feedback modes and stage 2).

\paragraph{What holds (equilibrium and elasticity).}
\begin{itemize}
  \item Electrostatics with same-sign images: exact shielding in circles, the depth law (0.984$\times$ prediction + 0.61 cells, $r=0.996$), the approach barrier,
        the weak leftover field in ellipses (Sec.~\ref{sec:theory}).
  \item Fluctuations of the passive wall lattice: along the wall, the exact potential with the measured bend constant $K_3=23\,k_BT$ (plus the detector's cell
        rounding) reproduces the rms within 2--4\%; across the wall a 0.15--0.5 cell excess remains, not explained by fluid noise (random walks); mobility
        $\mu\approx0.02$ (Sec.~\ref{sec:fluct}).
  \item Tilt and position are one variable: for a wall $+\half$, tail angle $-$ normal $=(2k-2)\phi+$const (lock 0.99 in $\approx$150\,000 active frames). The $2k-2$
        sites are where the tails point straight at the wall (heads into the bulk).
  \item Splay $\neq$ bend ($K_1\approx1.8K_3$) pins the active lattice at exactly those sites (barrier 80--110\,$k_BT$ at $k=6$; predicted collective rms 0.5 cells,
        measured 0.3). The \emph{passive} registry is not pinned: runs started from different initial directors settle at different registries and only creep; the
        shared registry of the production runs came from their common initial condition (B.~Loewe's spontaneous-symmetry-breaking hypothesis, confirmed).
        The \emph{active} registry is pinned: eight runs from six different initial directors all sit at the predicted site (after removing the tilt with the
        lock, $D=0.03\pm0.07$, resultant length 1.00); only the sense of rotation is chosen spontaneously.
\end{itemize}

\paragraph{Why the sites are fixed although the cavity is round (symmetry argument).}
\begin{itemize}
  \item \textbf{The winding breaks rotational symmetry to a discrete one.} The code imposes the wall director at lab angle $\theta=k\phi$ (\texttt{lc.c}, $O=0$),
        i.e.\ at $(k-1)\phi$ from the wall normal: homeotropic where $(k-1)\phi\equiv0$ (mod $\pi$), tangential halfway between. A rotation of the cavity by
        $\beta$ turns the pattern into $k\phi-(k-1)\beta$, the same pattern only if $(k-1)\beta\equiv0$ (mod $\pi$): the boundary has a $(2k-2)$-fold symmetry,
        with homeotropic points at $\phi_n=2\pi n/(2k-2)$. Only $k=1$ is fully rotationally symmetric.
  \item \textbf{The active sites are the homeotropic points.} The measured registry (phase 0) puts the $2k-2$ wall $+\half$ at $\phi_n$ with their tails
        (the splay side, where the director is radial) pointing straight at the wall: the radial side of the comet sits where the anchoring is radial. With a
        tilt $\tau$ the lattice is displaced from these points by $\tau/(2k-2)$, in the direction set by the sense of rotation.
  \item \textbf{One-constant elasticity hides this by an extra symmetry.} With a single Frank constant a uniform rotation of the director, independent of space,
        costs nothing; rotating the defects by $\beta$ and adding $(k-1)\beta$ to the director restores the boundary condition at no cost, hence the free sliding mode
        of Sec.~\ref{sec:fluct}. Splay $\neq$ bend breaks this internal symmetry (a director rotated relative to space changes the splay--bend mix) and turns the
        homeotropic points into energy minima.
  \item \textbf{Why the passive ring of $2k$ cannot be untilted.} By the tilt--position lock, a wall $+\half$ at $\phi_j$ has tilt (tail angle minus normal)
        $\tau_j=2\psi_\mathrm{rel}(\phi_j)=2(k-1)\phi_j+\text{const}$ (mod $2\pi$): its tail points at the wall only on a homeotropic point. All tilts vanish only if
        every defect sits on one of the $2k-2$ homeotropic points, impossible for $2k$ defects. For $N$ evenly spaced defects consecutive tilts differ by
        $2(k-1)\cdot2\pi/N$ (mod $2\pi$), zero only if $N$ divides $2(k-1)$; for $N=2k$ the step is $-2\pi/k$, so the tilts wind twice ($-4\pi$) around the ring.
        Measured ($k=6$): $+41^\circ,-20^\circ,-91^\circ,\dots$, steps of $-60^\circ=-2\pi/6$ (Sec.~\ref{sec:fluct}). The passive ring is therefore frustrated:
        most of its defects are tilted whatever the registry, which is why the registry energy averages out (incommensurate, $<1\,k_BT$).
        \emph{Interpretation:} the active state ejects exactly two $+\half$ into the bulk; the remaining $2k-2$ are the largest commensurate ring, which can
        sit untilted on the homeotropic points, so the $+1$ bulk also relieves the tilt frustration of the wall.
  \item \textbf{Commensurability decides how strongly.} The active lattice ($2k-2$ defects) matches the $(2k-2)$-fold pattern: barrier 80--110\,$k_BT$ ($k=6$),
        pinned at the predicted site from every initial condition, only the sense of rotation chosen spontaneously. The passive lattice ($2k$ defects) is
        incommensurate: the lowest harmonics cancel, the barrier is $<1\,k_BT$, and its registry is frozen in from the initial condition and only creeps.
  \item \textbf{Tested:} with a sine-modulated winding $k\phi+A\sin(m\phi)$ (branch \texttt{anchmod}) the active sites follow the unevenly spaced homeotropic points
        (0.21--0.27 cells, any initial director). The passive registry stays free: a modulation with $m$ not a multiple of $2k$ does not couple to an evenly spaced
        $2k$ ring at first order (Sec.~\ref{sec:sitecurv}).
\end{itemize}

\paragraph{What changed (hydrodynamics).}
\begin{itemize}
  \item \textbf{The wall layer does not drive the circulation; it brakes it.} Wall $+\half$ lean with their tails downstream, i.e.\ their heads point into the bulk and
        upstream; their active forces oppose the flow (Sec.~\ref{sec:hydro}, correction box). Withdrawn: ``the wall layer explains 70\% of the flow'', the force
        estimates $f/\eta$ per wall defect, ``the wall layer sustains the circulation'', and the reading of stage 2 as the layer selecting one vortex.
  \item \textbf{The bulk active stress drives the flow, including the two vortices of elongated ellipses.} Stokes flow of the full measured active force
        $-\nabla\!\cdot Q$ (no defect idealisation): one positive coefficient explains 48--72\% of the measured flow; the bulk part alone 47--73\%; the wall band
        opposes the circulation; the bulk reproduces the two-vortex part in ellipses (correlation 0.87--0.91) (Sec.~\ref{sec:dirstokes}).
  \item \textbf{Point-defect hydrodynamics is not enough.} Bulk $+\half$ as point forces explain only 6--21\%; a particle Langevin model with Stokeslets and log
        elasticity neither keeps $2k-2$ defects at the wall nor makes a circulation (Sec.~\ref{sec:defectmodel}).
\end{itemize}

\paragraph{What is new (the layer as an effective boundary).}
\begin{itemize}
  \item Seen from the bulk, the layer converts the $k$-fold patterned anchoring into an effectively \emph{planar} winding-1 boundary (exact, one-constant), smooth
        from $\approx0.34\ell$ from the wall. A tilt $\tau$ makes it a spiral, $c=90^\circ+\tau/2$, whose extensile stress drives the same circulation that tilts the
        layer (Sec.~\ref{sec:effbc}).
  \item This closes a feedback loop through the elastic field, not through forces at the wall. Its linear stability, with calibrated parameters and one assumed
        core size, gives an onset of tilt and circulation at $R/\ell_a\approx0.4$--0.6, where it is measured (Sec.~\ref{sec:effstab}); the predicted $k$-ordering is
        untested.
  \item Spacing does not give an independent criterion for the $+1$ bulk within the sweep: its existence is set by size ($R/\ell_a\ge0.76$, 99\%); $k=2$
        ($N=2$, smoothing depth $\approx R$) is predicted to be anomalous (one data point agrees).
\end{itemize}

\paragraph{Particle model with the spiral flow.} Given the observed structure (layer of $2k-2$ held at the wall, bulk pair kept inside), adding the
spiral flow to the defect Langevin model produces a circulation of the measured size and a tilt that follows it, above $R/\ell_a\approx0.7$ ($k=6$) and
later for $k=4$; without it, nothing. The model still cannot produce the structure itself (Sec.~\ref{sec:effstab}).

\paragraph{The size of the tilt.} To first order, splay $\neq$ bend only resists the tilt (exactly $\propto\cos\tau$). Beyond first order it adds a
$\cos2\tau$ term of the right sign ($+5\,k_BT$ at $k=6$), which makes the onset supercritical with a locked finite tilt; in the particle model this
locks the tilt at $\approx35^\circ$ with the computed size and at $63^\circ$ with a 3--4$\times$ larger term (Sec.~\ref{sec:effstab}).

\paragraph{Flow alignment.} Present in MPCD ($\lambda=2$, Leslie angle $30^\circ$), absent from the models. \emph{Corrected 2026-10-09:} the director at the wall defects' depth is rotated by 26--33$^\circ$
(Sec.~\ref{sec:tiltangle}; the earlier 4--5$^\circ$ estimate was wrong); it but erases the elastic spiral in the bulk (measured $74$--$87^\circ$ against the elastic $122^\circ$); the azimuthal active drive is confined to a ring
just inside the layer ($r/R\approx0.6$--0.9), braked in the wall band (Sec.~\ref{sec:effstab}).

\paragraph{Wall defects are not advected.} Measured: the wall $+\half$ do not move with the fluid at their depth (advection coefficient 0.002--0.004 against
a fluid speed of 0.13--0.16): the tilted layer sits at a static shifted registry. A particle model that advects them with the fluid slides instead of locking, even
with the measured drive (Sec.~\ref{sec:effstab}); the defect kinematics at an anchored wall is the next ingredient.

\paragraph{Site arrangement and curvature.} Moving the active sites with a modulated winding keeps them on the predicted homeotropic points (0.21--0.27 cells,
independent of the initial director) with the tilt unchanged; the passive ring is not pinned by a low-$m$ modulation (it couples only through $\sum e^{im\phi_j}=0$). In ellipses the layer is $\approx14\times$ leakier at $\varepsilon\ge1.5$ in one- and two-vortex intervals alike, with escapes at
the tips: the second vortex is not the immediate cause. A circle sweep with uneven sites (up to $A=3$) leaves every property of the layer unchanged: the site arrangement
does not destabilise it, so the ellipse instability comes from the curvature or the elongated flow (Sec.~\ref{sec:sitecurv}).

\paragraph{Imposed vortex.} With an imposed swirl the active layer behaves as a stick--slip system: pinned until the flow at its depth reverses, then a jump
of the size predicted by the lock. A passive layer under imposed flow also tilts to $\approx\pm60^\circ$ and slides with it. In the cavities the tilt is an active effect (the flow is active), but
its angle is set by the layer's response to flow: it grows with the flow-alignment parameter ($\approx32^\circ$ at $\lambda=0.5$, 60$^\circ$ at 1.2--2, 79$^\circ$ at 4),
not by the simple $2\theta_L$ rule (Sec.~\ref{sec:swirl}).

\paragraph{Half-integer windings.} For $k\ge2.5$ the odd layers (3, 5, 9 wall $+\half$) are pinned on the odd set of homeotropic points with the same lock,
tilt and single vortex as integer $k$; at $k=1.5$ (one site) no layer forms and the flow is weak and erratic, as the effective-boundary picture requires at least two sites;
$k=0.5$ gives a single bulk $+\half$ with a strong vortex (Sec.~\ref{sec:halfk}).

\paragraph{Why the active wall holds $2k-2$: flow-induced ejection.} (B.~Loewe's mechanism.) Flow aligns the wall defects to a common tilt; a $2k$ ring cannot take a
common tilt (its tilts must wind), which costs 375--475\,$k_BT$ at $k=6$ (200--510 for $k=4$--8), about twice the 180--200\,$k_BT$ (155--200) of ejecting two defects; the ejected $2k-2$ ring tilts by a registry shift
almost for free. Directly confirmed: \emph{passive} cavities under an imposed swirl eject two defects (12$\to$10 at the wall, bulk $+1$) for $\lambda=2$, 4, partly at
$\lambda=1.2$, never for tumbling $\lambda=0.5$. Activity enters by providing the flow (Sec.~\ref{sec:ejection}).

\paragraph{Ejection threshold and negative layer.} Passive ejection needs a flow strong enough to impose a common tilt (constant swirl: none up to 0.003, both
seeds at 0.01, tilt coherence 0.2--0.5 $\to$ 0.9). The mirror prediction for $-k$ (the wall gains two $-\half$) is favoured energetically but did not happen in passive runs:
it needs pair creation, rare without activity (Sec.~\ref{sec:ejection2}).

\paragraph{Shapes.} A single transition at aspect 1.25--1.37, scale free and the same for ellipses and ovals: coherence, circulation and the single vortex collapse together,
and the layer then loses defects in proportion to the size. Coherent flow imposed on a \emph{passive} ellipse gives the circle structure. In \emph{active} ellipses the imposed vortex
did not repair the layer, but it also never made the flow circle-like coherent, so activity disrupting the coherence remains a viable cause; stronger forcing is being tested
(Sec.~\ref{sec:shapesweep}).

\paragraph{Open, in order.} (1) Circulation of small high-$k$ cavities to
test the predicted $k$-ordering of the onset. (2) Effective-boundary picture in ellipses: does the bulk active stress of a spiral boundary of an ellipse split
into two vortices, and at which $\varepsilon$? (3) Rewrite of the hydrodynamics subsection around the director-field model (the old text is kept under the
correction box for traceability). (4) What holds the $+1$ in the bulk and the layer at the wall in the active state (the particle model needs it as input). 
